\documentclass[11pt,a4paper]{article}

\usepackage[utf8]{inputenc}
\usepackage[T1]{fontenc}
\usepackage[brazilian]{babel}

\usepackage{amsmath, amsthm, amssymb}

\usepackage{listings}
\usepackage{xcolor}
\usepackage{hyperref}
\usepackage{geometry}
\usepackage{enumitem}

\geometry{margin=2.5cm}

\hypersetup{
    colorlinks=true,
    linkcolor=blue,
    urlcolor=blue,
    citecolor=blue,
    pdftitle={Relatório Técnico - INF05027 Trabalho Final},
    pdfauthor={André Schaidhauer Luckmann}
}

\newtheorem{lemma}{Lema}
\newtheorem{theorem}{Teorema}

\lstset{
    basicstyle=\ttfamily\small,
    keywordstyle=\color{blue},
    commentstyle=\color{gray}\itshape,
    stringstyle=\color{red!70!black},
    showstringspaces=false,
    breaklines=true,
    frame=single,
    numbers=left,
    numberstyle=\tiny\color{gray},
    language=Python
}

\title{Relatório Trabalho Final \\ \large INF05027 - Projeto e Análise de Algoritmos I \\ Turma B \\ Trabalho Final}
\author{André Schaidhauer Luckmann \\ Cartão UFRGS: \texttt{00601117}}
\date{Junho de 2026}

\begin{document}

\maketitle

\vspace{1em}

\tableofcontents

\vspace{2em}

\section{Introdução}

Este relatório descreve as soluções desenvolvidas para o trabalho final da
disciplina de \textbf{Projeto e Análise de Algoritmos I}, cujo
objetivo é consolidar os estudos de algoritmos sobre grafos por meio da
implementação e análise de soluções para dois problemas na plataforma Beecrowd:

\begin{itemize}
    \item \textbf{Problema 1:} \emph{O Rato no Labirinto}
    (\url{https://judge.beecrowd.com/pt/problems/view/1799}).
    \item \textbf{Problema 2:} \emph{Ir e Vir}
    (\url{https://judge.beecrowd.com/pt/problems/view/1128}).
\end{itemize}

A linguagem escolhida para ambas as implementações foi \textbf{Python}
por já ser utilizada como linguagem principal da disciplina, possuindo vários templates de código que estudamos e por ser uma linguagem intuitiva que oferece bibliotecas com diversas estruturas de dados completas. Além disso, o autor (André) já possuía experiência e familiaridade com \textbf{Python}.

\section{Problema 1: O Rato no Labirinto}

\subsection{Descrição do Problema}

O ratinho IBO deve atravessar um
labirinto representado como um grafo não-direcionado. Ele parte de um vértice
rotulado \texttt{Entrada}, deve obrigatoriamente passar pelo vértice que
contém o \texttt{queijo}, identificado pelo caractere \texttt{*}, e
terminar no vértice \texttt{Saida}. IBO deve seguir
sempre o caminho mais curto em seus trajetos. Caso seja necessário passar por um mesmo
ponto duas vezes (uma indo até o queijo, outra voltando até a saída), \emph{cada
visita é contada separadamente}.

Dada a descrição do labirinto na entrada (número de pontos, número de
ligações e a lista de pares de vértices ligados), o objetivo é imprimir
o número mínimo de pontos pelos quais IBO passa para cumprir a tarefa.

\subsection{Estratégia da Solução}

\paragraph{Modelagem.} O labirinto é representado como um grafo
não-direcionado $G=(V,E)$, em que $V$ é o conjunto de pontos rotulados
do labirinto (palavras compostas apenas por letras, como \texttt{A},
\texttt{F}, \texttt{Entrada}, \texttt{Saida}, e o queijo \texttt{*})
e $E$ é o conjunto de ligações. As arestas do grafo não são valoradas, pois não é fornecida essa informação e estamos interessados na quantidade de pontos visitados (tamanho do caminho).

\paragraph{Decomposição.} A observação central é que o caminho ótimo
\texttt{Entrada} $\to \cdots \to *\to \cdots \to$ \texttt{Saida} pode
ser decomposto em dois caminhos mínimos \emph{independentes}:

\begin{enumerate}[label=(\alph*)]
    \item caminho mínimo de \texttt{Entrada} até \texttt{*};
    \item caminho mínimo de \texttt{*} até \texttt{Saida}.
\end{enumerate}

Como o problema explicita que cada visita conta separadamente, não há
ganho em tentar reaproveitar trechos comuns: a soma dos comprimentos
dos dois caminhos mínimos é justamente o número total mínimo de
visitas, conforme demonstrado na Seção~\ref{sec:correcao1}.

\paragraph{Algoritmo.} Aplica-se Busca em Largura (BFS) duas vezes:

\begin{enumerate}
    \item \texttt{bfs\_distance(graph, start\_node="Entrada", target\_node="*")}
    \item \texttt{bfs\_distance(graph, start\_node="*", target\_node="Saida")}
\end{enumerate}

\noindent
e a resposta é a soma dos dois valores retornados.

\paragraph{Estruturas de dados.}
\begin{itemize}
    \item \texttt{graph: Dict[Node, List[Node]]} - grafo representado como lista de adjacências
    usando um dicionário.
    \item \texttt{collections.deque} - Double-ended queue, utilizado somente como fila de vértices a serem visitados.
    \item \texttt{visited: Dict[Node, bool]} - controle de vértices
    já visitados.
    \item \texttt{distances: Dict[Node, int]} - distância (em número de arestas)
    de cada vértice ao vértice de partida da busca em largura.
\end{itemize}

\paragraph{Justificativa de design.} A BFS é \emph{ótima} para encontrar
caminhos mínimos em grafos não valorados, com custo linear no tamanho
do grafo. O algoritmo de Dijkstra foi considerado inicialmente, mas teria um custo
adicional desnecessário (heap). Como a busca em largura vai passando por "níveis" do grafo, ela é uma forma mais otimizada que a DFS caso o algoritmo de DFS entre em uma ramificação muito grande do grafo que não leve ao nodo desejado (pior caso).

\subsection{Argumento de Correção}
\label{sec:correcao1}

\begin{lemma}[Corretude da BFS em grafos não valorados]
\label{lem:bfs}
A BFS, partindo de um vértice $s$, atribui a cada vértice $u$ alcançável
exatamente a distância em arestas $d(s,u)$.
\end{lemma}

\begin{proof}[Esboço]
Por indução sobre as camadas de descoberta. A camada $0$ é $\{s\}$,
com $d=0$. Suponha que ao final do processamento da camada $k$ todos
os vértices alcançáveis a distância exatamente $k$ tenham sido marcados
e enfileirados com distância correta. Quando um vértice $v$ da camada
$k$ é desenfileirado, examina-se cada vizinho $w \in N(v)$: se $w$
ainda não está marcado, ele é colocado na fila com
$\texttt{distances}[w] = \texttt{distances}[v] + 1 = k+1$. Como a fila
obedece a lógica FIFO (First-In-First-Out) e somente vértices das camadas anteriores foram enfileirados
antes, nenhum vértice cuja distância seja menor que $k+1$ poderia
ainda estar não marcado. Logo, todos os vértices à distância $k+1$ são
descobertos com a distância correta. Pela indução, a BFS é correta.
\end{proof}

\begin{lemma}[Decomposição]
\label{lem:decomp}
Seja $\mathrm{Opt}$ o número mínimo de pontos visitados em qualquer
caminho \texttt{Entrada} $\to \cdots \to *\to \cdots \to$ \texttt{Saida}
em $G$. Então
\[
    \mathrm{Opt} \;=\; d(\texttt{Entrada}, *) \;+\; d(*, \texttt{Saida}),
\]
onde $d(\cdot,\cdot)$ é a distância em arestas em $G$.
\end{lemma}

\begin{proof}[Esboço]
\textbf{($\le$)} Sejam $P_1$ um caminho mínimo de \texttt{Entrada} a
\texttt{*} (com $d(\texttt{Entrada},*)$ arestas, ou seja,
$d(\texttt{Entrada},*) + 1$ vértices contados sem repetição), e $P_2$
um caminho mínimo análogo de $*$ a \texttt{Saida}. A concatenação
$P_1 \cdot P_2$ é um caminho válido para o problema, e como o
enunciado conta a passagem por \texttt{*} duas vezes (uma vez ao final
de $P_1$ e outra ao início de $P_2$), o número total de visitas é
exatamente $|P_1| + |P_2| = d(\texttt{Entrada},*) + d(*,\texttt{Saida})$.
Logo, $\mathrm{Opt} \le d(\texttt{Entrada},*) + d(*,\texttt{Saida})$.

\textbf{($\ge$)} Reciprocamente, qualquer caminho válido $P$ pode ser
particionado pela primeira ocorrência de \texttt{*} em sub-caminhos
$P_1'$ (de \texttt{Entrada} a \texttt{*}) e $P_2'$ (de \texttt{*} a
\texttt{Saida}). Cada um deles é um caminho em $G$, portanto
$|P_1'| \ge d(\texttt{Entrada},*)$ e $|P_2'| \ge d(*,\texttt{Saida})$.
Somando, o número total de visitas em $P$ é
$\ge d(\texttt{Entrada},*) + d(*,\texttt{Saida})$. Em particular,
$\mathrm{Opt} \ge d(\texttt{Entrada},*) + d(*,\texttt{Saida})$.

Combinando as duas desigualdades, vale a igualdade.
\end{proof}

\paragraph{Validade da implementação.} As duas chamadas de
\texttt{bfs\_distance} terminam, pois o grafo é finito e cada vértice
é enfileirado no máximo uma vez (após a marcação). O resultado é
sempre não-negativo. Caso o vértice-alvo seja inalcançável, a função
retorna $-1$; o enunciado garante a existência dos caminhos, de modo
que esse caso não ocorre nas instâncias do problema.

\subsection{Análise de Complexidade}

Dados os limites de Pontos ($4 \le Pontos \le 4000$) e Ligacoes ($4 \le Ligacoes \le 5000$), sejam $n = |V| \le 4000$ e $m = |E| \le 5000$. Analisa-se função a
função:

\begin{itemize}
    \item \texttt{build\_graph\_from\_input}: lê $m$ linhas. Cada
    inserção em \texttt{dict} e em \texttt{list} é $O(1)$ amortizado.
    \textbf{Tempo:} $O(m)$. \textbf{Memória:} $O(n+m)$ (lista de
    adjacências).
    \item \texttt{bfs\_distance}: a inicialização de \texttt{visited}
    e \texttt{distances} percorre todos os vértices, custando $O(n)$.
    O laço principal enfileira cada vértice no máximo uma vez ($O(n)$)
    e percorre cada aresta no máximo duas vezes (uma por extremidade), custando $O(m)$. \textbf{Tempo
    total:} $O(n+m)$. \textbf{Memória adicional:} $O(n)$.
    \item \texttt{main}: realiza duas chamadas independentes de BFS.
    Custo total: $2 \cdot O(n+m) = O(n+m)$.
\end{itemize}

\paragraph{Resumo.} \textbf{Complexidade de tempo:} $O(n+m)$.
\textbf{Complexidade de memória:} $O(n+m)$, dominada pela lista de
adjacências.

\subsection{Discussão}

\begin{itemize}
    \item \textbf{Dificuldades encontradas.} O ponto não-trivial foi
    interpretar corretamente a regra de contagem de visitas: como cada
    passagem por um ponto conta como uma visita independente, a
    decomposição em duas BFS independentes torna-se ótima e simples.
    Sem essa observação, é tentador buscar uma solução mais elaborada
    que tente compartilhar trechos.

    \item \textbf{Alternativas consideradas.}
    \begin{itemize}
        \item \emph{Dijkstra:} aplicável, mas com sobrecarga de fila
        de prioridade desnecessária para arestas unitárias.
        \item \emph{DFS:} não garante caminho mínimo, portanto
        inadequada.
    \end{itemize}

    \item \textbf{Possíveis otimizações.}
    \begin{itemize}
        \item Encerrar a BFS no instante em que o vértice-alvo é
        \emph{desenfileirado}, em vez de detectá-lo na inserção, é
        uma refatoração equivalente em corretude e talvez mais
        idiomática.
    \end{itemize}

    \item \textbf{Aprendizados.} A BFS pode ser vista como caso
    particular de Dijkstra quando todas as arestas têm o mesmo peso.
\end{itemize}

\section{Problema 2: Ir e Vir}

\subsection{Descrição do Problema}

Em uma certa cidade, há $N$ interseções ligadas por ruas que podem ser
de mão única ou mão dupla. Deseja-se verificar a propriedade de
\emph{conexidade}: dadas quaisquer duas interseções $V$ e $W$, é
possível viajar de $V$ para $W$ \emph{e} de $W$ para $V$.

A entrada contém vários casos de teste que fornecem $N$, $M$, seguidos de
$M$ linhas descrevendo $M$ ruas no formato $(V, W, P)$ onde $P=1$ indica via de mão única
$V \to W$ e $P=2$ indica via de mão dupla $V \leftrightarrow W$. A
sequência termina com uma linha contendo \texttt{N=0 e M=0}. Para cada
caso, deve-se imprimir $1$ se a propriedade é satisfeita ou $0$ caso
contrário.

\subsection{Estratégia da Solução}

\paragraph{Modelagem.} Cada caso de entrada é modelado como um dígrafo
$G=(V,A)$. Uma rua de mão única $V \to W$ contribui com a aresta
$(v,w)$ e uma rua de mão dupla $V \leftrightarrow W$ contribui com as duas arestas $(v,w)$ e $(w,v)$.
A propriedade do enunciado ``ir de $V$ a $W$ e voltar de $W$ a
$V$ para todo par'' é uma característica de 
um componente fortemente conexo em um grafo. Desta forma, para resolver o problema, a interpretação foi de que deve-se verificar se $G$ é composto por um
único componente fortemente conexo (SCC, Strongly-Connected-Component) que contém todos os
vértices.

\paragraph{Algoritmo.} Aplica-se o algoritmo
\textbf{Kosaraju-Sharir} para encontrar os SCCs, com uma otimização:
basta saber se há \emph{exatamente um} SCC. Portanto, ao detectar a
necessidade de iniciar uma segunda busca DFS, o
algoritmo retorna $0$ imediatamente.

\paragraph{Etapas.}
\begin{enumerate}
    \item Construir o grafo invertido $G^T$
    (\texttt{reverse\_graph}).
    \item Executar uma DFS em $G^T$ produzindo a ordem reversa de
    finalização (ordenação topológica via DFS recursivo, função
    \texttt{dfs\_toposort}).
    \item Executar uma DFS em $G$ visitando os vértices na ordem
    obtida na etapa anterior pela função \texttt{dfs\_toposort}. Cada chamada de DFS a partir de um vértice ainda não visitado descobre exatamente um SCC.
    \item Contar quantas chamadas de DFS são necessárias para cobrir
    todos os vértices: se for mais que uma, retornar $0$; caso
    contrário, retornar $1$.
\end{enumerate}

\paragraph{Estruturas de dados.}
\begin{itemize}
    \item \texttt{graph: Dict[Node, List[Node]]} - grafo representado como lista de adjacências
    usando um dicionário.
    \item \texttt{rev\_g: Dict[Node, List[Node]]} - grafo inverso representado como lista de adjacências
    usando um dicionário.
    \item \texttt{visited: Dict[Node, bool]} - controle de vértices
    já visitados.
    \item \texttt{order: list[Node]} - lista com a ordem inversa de vértices visitados por DFS
    em $G^T$.
\end{itemize}

\subsection{Argumento de Correção}

\begin{theorem}[Kosaraju-Sharir]
\label{teo:kosaraju}
Os componentes fortemente conexos de um dígrafo $G$ são as árvores da floresta DFS executada em $G$ visitando-se os vértices
na ordem inversa de visita da DFS feita em $G^T$.
\end{theorem}

\begin{proof}[Esboço]
Seja $\mathrm{SCC}(G)$ o conjunto de SCCs de $G$ e considere o
\emph{grafo de componentes} $\mathcal{C}(G)$, que é acíclico (DAG). A primeira DFS em
$G^T$ visita os nodos em uma ordem determinada. A propriedade chave é: o último vértice
a ser visitado na DFS em $G^T$ pertence a um
SCC que é \emph{fonte} em $\mathcal{C}(G^T)$, e ao mesmo tempo é
\emph{sumidouro} em $\mathcal{C}(G)$. Iniciando uma DFS em $G$ a
partir desse vértice, descobrem-se exatamente os vértices do seu
SCC, pois \emph{nenhuma aresta sai} desse SCC em $G$ para outros
SCCs (por ele ser sumidouro em $\mathcal{C}(G)$). Por indução,
removendo o SCC já descoberto, o argumento se repete para o próximo
vértice de maior tempo de finalização ainda não visitado. Logo, cada
chamada de DFS na fase final descobre exatamente um SCC.
\end{proof}

\paragraph{Adaptação à implementação.} O código aplica a
ordenação topológica por DFS no grafo \emph{invertido}
(\texttt{rev\_g}) e a DFS final no grafo \emph{original}
(\texttt{graph}). Como $(G^T)^T = G$, a ordem produzida é
equivalente à ordem inversa de visitação exigida
pelo Teorema~\ref{teo:kosaraju}. Em particular, cada execução de
\texttt{dfs(graph, root)} dentro do laço final descobre um único
SCC de $G$.

\paragraph{Critério de aceitação.} O dígrafo é fortemente conexo
se, e somente se, $|\mathrm{SCC}(G)| = 1$, isto é, uma única
chamada de DFS na fase final cobre todos os vértices. Na função
\texttt{is\_connected}, o contador \texttt{connected\_components}
é incrementado a cada nova raiz de DFS; quando uma segunda raiz
seria iniciada, a função retorna $0$.

\paragraph{Validade.} O grafo é finito e ambas as DFS terminam.
A saída é sempre binária. O \emph{early-exit} não compromete a
corretude, pois bastam dois SCCs distintos para que a resposta seja
$0$.

\subsection{Análise de Complexidade}

Sejam $n = |V| \le 2000$ e $m = |A| \le n(n-1)/2$. Por caso de
teste:

\begin{itemize}
    \item \texttt{build\_graph\_from\_input}: São lidas $m$ linhas, logo a complexidade de tempo é $O(m)$. Como o grafo é armazenado em memória como uma lista de adjacência usando um dicionário, o custo de memória é $O(n+m)$.
    \item \texttt{reverse\_graph}: percorre cada par de (nodo, vizinho) uma vez ao construir as adjacências invertidas. $O(n+m)$ tempo
    e $O(n+m)$ memória.
    \item \texttt{dfs\_toposort}: a DFS clássica visita cada vértice
    uma vez (custo $O(n)$) e cada aresta uma vez (custo $O(m)$);
    a inversão de \texttt{order} ao final é $O(n)$. Complexidade de tempo total é $O(n+m)$. 
    O dicionário \texttt{visited} requer memória adicional com custo $O(n)$.
    \item \texttt{is\_connected}: além das chamadas acima, executa
    no máximo uma DFS completa adicional. Pior caso $O(n+m)$.
\end{itemize}

\paragraph{Resumo por caso de teste:} \textbf{Tempo:} $O(n+m)$,
\textbf{Memória:} $O(n+m)$.

\paragraph{Custo total.} Para $t$ casos de teste, o custo total é
$O\bigl(t \cdot (n+m)\bigr)$.

\subsection{Discussão}

\begin{itemize}
    \item \textbf{Decisões.} A percepção mais
    importante foi reconhecer que a propriedade ``viajar de $V$ a
    $W$ e de $W$ a $V$ para todo par'' é uma característica de um componente fortemente conexo em um dígrafo. A partir disso, consultando os materiais de aula e \href{https://github.com/BrunoGrisci/projeto-e-analise-de-algoritmos/blob/main/paa1/graph_search.ipynb}{repositório no Github}, foi encontrado o algoritmo de Kosaraju-Sharir que encontra os SCCs de um dígrafo. Dessa forma, analisando o problema foi constatado que o programa deveria verificar se o dígrafo é um SCC que contém todos os nodos do dígrafo.

    \item \textbf{Alternativas consideradas.} Inicialmente, havia pensado em realizar buscas DFS/BFS entre um nodo fonte e um nodo alvo, para todos os nodos do grafo, e caso não fosse possível satisfazer algum caminho, o programa deveria retornar $0$. Porém, essa alternativa provavelmente seria força bruta e ineficiente, resultando em uma complexidade de ordem cúbica $O(n^3)$.

    \item \textbf{Possíveis otimizações.}
    \begin{itemize}
        \item Como o algoritmo apenas precisa saber se há um único
        componente fortemente conexo, o contador 
        \texttt{connected\_components} poderia ser
        substituído por uma simples \texttt{flag} booleana.
    \end{itemize}

    \item \textbf{Aprendizados.} O algoritmo de Kosaraju-Sharir mostra como
    a ordem de visitas da DFS no grafo transposto induz uma
    ordem topológica do DAG de componentes. Também demonstra a
    propriedade dos sumidouros do grafo de componentes: o SCC do
    último vértice visitado em $G^T$ é um sumidouro
    em $\mathcal{C}(G)$, do qual nenhuma aresta sai para outros
    SCCs.
\end{itemize}

\section{Conclusão}

Este trabalho consolidou conceitos centrais de algoritmos sobre
grafos por meio da implementação de duas soluções para problemas
de Grafos:

\begin{itemize}
    \item Para o problema \textbf{O Rato no Labirinto}, foi mostrado que o problema
    de caminho mínimo passando obrigatoriamente por um vértice
    intermediário pode ser decomposto em
    duas BFSs independentes em um grafo não valorado, resultando em uma solução
    linear no tamanho do grafo, com custo $O(n+m)$.

    \item Para o problema \textbf{Ir e Vir}, foi aplicado o algoritmo
    de Kosaraju-Sharir para verificar a quantidade de componentes fortemente conexos do
    dígrafo, retornando $0$ em caso de detecção de mais de um componente fortemente conexo no dígrafo. A solução possui complexidade $O(n+m)$ por caso de teste.
\end{itemize}

Ambas as soluções respeitam os limites de entrada do
juiz e foram aceitas pelo Beecrowd. Além do desenvolvimento de código correto, o exercício reforçou a importância de modelar adequadamente o problema (identificar se o grafo é direcionado ou não, peso das arestas, decomposição em sub-problemas) e de fornecer
argumentos de correção formais e análise de complexidade.

\end{document}
